\chapter{Условно математическо очакване}

\section{Мотивация: най-добра прогноза}

\begin{rembox}[Отправна точка]
Знаем, че $DX = \mathbb E(X-\mathbb EX)^2 = \min_{a\in\mathbb R}
\mathbb E(X-a)^2$ - т.е.\ константата $a=\mathbb EX$ е
\emph{най-добрата} (в~смисъл на най-малка средноквадратична грешка)
константна прогноза за случайната величина $X$. Естественият
въпрос е: ако разполагаме с~\emph{допълнителна информация} за друга
случайна величина $Y$, каква е най-добрата функция $G(Y)$, която
прогнозира $X$, т.е.\ коя $G$ минимизира $\mathbb E\big(X-G(Y)\big)^2$?
\end{rembox}

\begin{rembox}[Извеждане за бинарно $Y$]
Нека $Y\in\{0,1\}$, $A=\{Y{=}1\}$, $A^c=\{Y{=}0\}$. Всяка функция на
$Y$ има вида $G(Y)=G(1)I_A+G(0)I_{A^c}=aI_A+
bI_{A^c}$ за някои константи $a,b$, така че задачата се
свежда до
\[
  f(a,b) := \mathbb E\big(X-aI_A-bI_{A^c}\big)^2 = \mathbb EX^2 + a^2\,\mathbb EI_A + b^2\,\mathbb EI_{A^c} - 2a\,\mathbb EXI_A - 2b\,\mathbb EXI_{A^c} \to \min_{a,b\in\mathbb R}.
\]
Полагаме частните производни на нула:
\begin{align*}
  0 &= \frac{\partial f}{\partial a} = 2a\,\mathbb EI_A - 2\,\mathbb EXI_A \ \Rightarrow\ a = \frac{\mathbb EXI_A}{\mathbb P(A)},\\
  0 &= \frac{\partial f}{\partial b} = 2b\,\mathbb EI_{A^c} - 2\,\mathbb EXI_{A^c} \ \Rightarrow\ b = \frac{\mathbb EXI_{A^c}}{\mathbb P(A^c)}
\end{align*}
(използвахме $\mathbb EI_A=\mathbb P(A)$). Следователно
оптималната прогноза е
\[
  G^*(Y) = \frac{\mathbb EXI_A}{\mathbb P(A)}\,I_A + \frac{\mathbb EXI_{A^c}}{\mathbb P(A^c)}\,I_{A^c},
\]
т.е.\ функцията, която най-добре предсказва $X$ на базата на $Y$, е
именно \textbf{условното математическо очакване}.
\end{rembox}

\section{Дефиниция}

\begin{defbox}[~(Условно математическо очакване)]
Нека $X,Y$ са случайни величини с~$\mathbb EX^2<\infty$ (т.е.\
$X\in L^2$, за да е добре дефиниран минимумът по-долу). Тогава
$\mathbb E(X\mid Y):=G^*(Y)$ е случайната величина, за която
\[
  \min_{G}\, \mathbb E\big(X-G(Y)\big)^2 = \mathbb E\big(X-G^*(Y)\big)^2 = \mathbb E\big(X-\mathbb E(X\mid Y)\big)^2,
\]
т.е.\ функцията на $Y$, която минимизира средноквадратичната грешка
на прогнозата за~$X$.
\end{defbox}

\begin{thmbox}[]{Теорема}
Нека $X,Y$ са случайни величини, като $Y$ е дискретна с~разпределение
$\mathbb P(Y=y_j)=p_j$. Полагаме $A_j=\{Y=y_j\}$, така че
$Y=\sum_j y_jI_{A_j}$. Тогава
\[
  \mathbb E[X\mid Y] = \sum_j \frac{\mathbb EXI_{A_j}}{\mathbb P(A_j)}\,I_{A_j}
\]
(т.е.\ $\mathbb E[X\mid Y]$ е случайна величина).
\end{thmbox}

\begin{proofbox}
Нека $G(Y)=\sum_j a_jI_{A_j}$ е произволна функция на $Y$
(параметризирана с~коефициентите $a_j$). Разглеждаме
\[
  f(a_1,a_2,\dots) = \mathbb E\big(X-G(Y)\big)^2 = \mathbb E\Big(X-\sum_j a_jI_{A_j}\Big)^2.
\]
За всяко $i$, полагаме частната производна по $a_i$ на нула (понеже
$\{A_j\}$ са несъвместими, кръстосаните членове изчезват):
\[
  0 = \frac{\partial f}{\partial a_i} = 2\,\mathbb E\Big(X-\sum_j a_jI_{A_j}\Big)I_{A_i} = 2\big(\mathbb EXI_{A_i} - a_i\,\mathbb EI_{A_i}\big)
  \ \Rightarrow\ a_i = \frac{\mathbb EXI_{A_i}}{\mathbb P(A_i)}.
\]
Следователно $\mathbb E(X\mid Y) = \sum_j \dfrac{\mathbb EX
I_{A_j}}{\mathbb P(A_j)}\,I_{A_j}$. \qedbg
\end{proofbox}

\begin{rembox}[Разписване по стойности на \texorpdfstring{$X$}{X}]
Нека допълнително $X=\sum_i x_iI_{B_i}$ е дискретна,
$B_i=\{X=x_i\}$. Тогава
\begin{align*}
  \frac{\mathbb EXI_{A_j}}{\mathbb P(A_j)} &= \frac{\mathbb E\sum_i x_iI_{B_i}I_{A_j}}{\mathbb P(A_j)} = \sum_i x_i\,\frac{\mathbb P(A_j\cap B_i)}{\mathbb P(A_j)}\\
  &= \sum_i x_i\,\mathbb P(B_i\mid A_j) = \sum_i x_i\,\mathbb P(X=x_i\mid Y=y_j),
\end{align*}
което мотивира следната по-фина дефиниция - стойността на условното
очакване \emph{при конкретно фиксирано} $Y=y_j$.
\end{rembox}

\begin{defbox}[~(УМО при фиксирано \texorpdfstring{$Y=y_j$}{Y=yj})]
Нека $X,Y$ са дискретни случайни величини. Тогава за всяко $j$:
\[
  \mathbb E[X\mid Y=y_j] = \sum_i x_i\,\mathbb P(X=x_i\mid Y=y_j).
\]
\end{defbox}

\begin{rembox}[Важно разграничение]
$\mathbb E[X\mid Y]$ е \textbf{случайна величина} (функция на
изхода на случайния експеримент чрез $Y$), докато $\mathbb E[X\mid
Y=y_j]$ е \textbf{число} - стойността на тази случайна величина,
когато $Y$ вземе конкретната стойност~$y_j$.
\end{rembox}

\begin{corbox}[]
Нека $X,Y$ са дискретни случайни величини. Тогава $\mathbb E[X\mid
Y]$ е дискретна случайна величина, приемаща стойностите $\mathbb
E[X\mid Y=y_j]$ с~вероятности $\mathbb P(Y=y_j)$:
\[
  \begin{array}{c|cccc}
    \mathbb E[X\mid Y] & \mathbb E[X\mid Y{=}y_1] & \dots & \mathbb E[X\mid Y{=}y_j] & \dots\\\hline
    \mathbb P & p_1=\mathbb P(A_1) & \dots & p_j=\mathbb P(A_j) & \dots
  \end{array}
\]
\end{corbox}

\section{Свойства на условното математическо очакване}

\begin{thmbox}[~(на условното математическо очакване)]{Свойства}
Нека $X,Y,Z$ са случайни величини с~$\mathbb EX<\infty$,
$\mathbb EZ<\infty$, като $Y$ е дискретна.
\begin{enumerate}[label=\arabic*)]
  \item \textbf{Линейност:} $\mathbb E[aX+bZ\mid Y] = a\,\mathbb E[X\mid Y]+b\,\mathbb E[Z\mid Y]$.
  \item \textbf{Независимост:} ако $X\perp\!\!\!\perp Y$, то $\mathbb E[X\mid Y] = \mathbb EX$.
  \item \textbf{„Изваждане на познатото“:} ако $X=f(Y)$ за някаква функция $f$, то $\mathbb E[X\mid Y]=X=f(Y)$.
  \item \textbf{Кула-свойство (закон за повторното очакване):} $\mathbb E\big[\mathbb E[X\mid Y]\big] = \mathbb EX$.
\end{enumerate}
\end{thmbox}

\begin{proofbox}
Навсякъде означаваме $A_j=\{Y=y_j\}$.

\textbf{1)}
\[
  \mathbb E[aX+bZ\mid Y] = \sum_j \frac{\mathbb E\big[(aX+bZ)I_{A_j}\big]}{\mathbb P(A_j)}\,I_{A_j} = \sum_j \frac{a\,\mathbb EXI_{A_j}+b\,\mathbb EZI_{A_j}}{\mathbb P(A_j)}\,I_{A_j}
\]
\[
  = a\sum_j\frac{\mathbb EXI_{A_j}}{\mathbb P(A_j)}I_{A_j} + b\sum_j\frac{\mathbb EZI_{A_j}}{\mathbb P(A_j)}I_{A_j} = a\,\mathbb E[X\mid Y]+b\,\mathbb E[Z\mid Y].
\]

\textbf{2)} Ако $X\perp\!\!\!\perp Y$, то $X\perp\!\!\!\perp\mathbb
1_{A_j}$ за всяко $j$ (индикаторът на $Y$ носи информация само за
$Y$), следователно $\mathbb EXI_{A_j} = \mathbb EX\cdot
\mathbb EI_{A_j} = \mathbb EX\cdot\mathbb P(A_j)$, т.е.
$\dfrac{\mathbb EXI_{A_j}}{\mathbb P(A_j)}=\mathbb EX$ за
всяко $j$. Тогава
\[
  \mathbb E[X\mid Y] = \sum_j\frac{\mathbb EXI_{A_j}}{\mathbb P(A_j)}I_{A_j} = \sum_j \mathbb EX\cdot I_{A_j} = \mathbb EX\sum_j I_{A_j} = \mathbb EX
\]
(понеже $\{A_j\}$ е пълна група от събития, $\sum_j I_{A_j}
\equiv1$).

\textbf{3)} Ако $X=f(Y)$, то върху $A_j=\{Y=y_j\}$ имаме
$X=f(y_j)$ (константа), следователно
\[
  \frac{\mathbb EXI_{A_j}}{\mathbb P(A_j)} = \frac{\mathbb E f(Y)I_{A_j}}{\mathbb P(A_j)} = \frac{f(y_j)\,\mathbb P(A_j)}{\mathbb P(A_j)} = f(y_j),
\]
и~оттам $\mathbb E[X\mid Y] = \sum_j f(y_j)I_{A_j} = f(Y) = X$.

\textbf{4)} Пряко от дефиницията и~линейността на очакването:
\begin{align*}
  \mathbb E\big[\mathbb E[X\mid Y]\big] &= \mathbb E\Big[\sum_j\frac{\mathbb EXI_{A_j}}{\mathbb P(A_j)}I_{A_j}\Big] = \sum_j\frac{\mathbb EXI_{A_j}}{\mathbb P(A_j)}\,\mathbb EI_{A_j}\\
  &= \sum_j \mathbb EXI_{A_j} = \mathbb EX\Big(\sum_jI_{A_j}\Big) = \mathbb EX,
\end{align*}
защото $\mathbb EI_{A_j}=\mathbb P(A_j)$ съкращава знаменателя,
а~$\{A_j\}$ е пълна група от събития. \qedbg
\end{proofbox}

\section{Условно разпределение}

\begin{defbox}[~(Условно разпределение)]
Нека $X,Y$ са дискретни случайни величини. \textbf{Условното
разпределение} на $X$ при положение $Y=y_j$ е
\[
  \begin{array}{c|ccc}
    X\mid Y{=}y_j & x_1 & \dots & x_i \dots\\\hline
    \mathbb P & \mathbb P(X{=}x_1\mid Y{=}y_j) & \dots & \mathbb P(X{=}x_i\mid Y{=}y_j)\dots
  \end{array}
\]
За всяко фиксирано $j$, $\sum_i \mathbb P(X=x_i\mid Y=y_j)=1$ - това
е коректно вероятностно разпределение (върху стойностите на $X$,
„при новите правила на играта“, зададени от условието $Y=y_j$).
\end{defbox}
